\chapter{Future Work}
\label{chap:future-work}

\section{Introduction}

The project completed its proposed experimental and implementation objectives,
but phishing methods and normal web services continue to change. Future work
should therefore concentrate on maintaining the detector, reducing avoidable
warnings, and evaluating the browser extension with broader evidence. The
following improvements extend the current system without changing the main
conclusion of this report.

\section{Newer and Broader URL Data}

The highest priority is to collect labelled URLs at different times after the
current dataset freeze. Later snapshots should contain verified phishing URLs,
ordinary browsing, long legitimate URLs, login and account pages, Arabic and
Libyan services, cloud documents, and user-controlled content hosted on trusted
platforms. Collection time and source should be stored so that later data can
be used as a true external test rather than mixed into the existing test sets.

When sufficient new data have been reviewed, a later model may be retrained
using a new training split. The present frozen tests should still be retained
for historical comparison, while an additional untouched snapshot becomes the
new final test. This would show whether an improvement transfers across time
instead of only fitting one evaluation period.

\section{Reducing False Warnings and Missed Phishing}

Future experiments can examine the errors with the highest practical cost.
Legitimate examples should include difficult login, redirect, documentation,
search, cloud-storage, and long-query URLs. Difficult phishing examples should
include compromised services, shared documents, short paths, and URLs whose
structure resembles a normal service.

Small controlled experiments with class weighting, focal loss, and hard-example
mining could test whether the score distributions become better separated.
Thresholds should continue to be chosen on validation data under stated FPR
constraints. The final test must remain untouched until the choice is frozen.
This preserves a fair comparison and avoids lowering FNR by simply accepting an
unusable number of false warnings.

\section{Additional Detection Evidence}

URL-only analysis is fast and privacy-friendly, but it cannot observe every
property of a website. A later system could combine the current URL model with
lightweight independent signals, including:

\begin{itemize}
\item domain and DNS history;
\item certificate and HTTPS information;
\item domain reputation and verified phishing feeds;
\item safe HTML or DOM features collected after explicit permission; and
\item redirect-chain and final-destination information.
\end{itemize}

These signals should be added as a secondary review layer rather than presented
as replacements for the GAN-augmentation experiment. A controlled comparison
of real-only and GAN-augmented primary models with the same secondary evidence
would show whether the augmentation benefit remains present.

\section{Further GAN Research}

Future generator work could use newer verified phishing data and compare
several generator checkpoints over time. Evaluation should retain the existing
validity, uniqueness, overlap, diversity, realism, and classifier-challenge
checks. Human review of a small defanged sample could be added with a clear
review protocol.

The selection rule should be fixed before classifier testing. Generated URLs
should not be selected because they happen to improve a final test result.
Experiments may also compare whether augmentation is more valuable when less
real phishing training data are available. Such work would clarify when
synthetic data supplies useful coverage and when additional real data is the
better choice.

\section{Browser Extension Improvements}

The browser prototype can be improved through:

\begin{itemize}
\item clearer explanations of caution, warning, and block decisions;
\item broader testing of Arabic localisation and right-to-left layouts;
\item accessibility checks for colour, keyboard use, and screen readers;
\item lower memory and initialization cost on older computers;
\item clearer management and expiry of user-allowed hostnames;
\item versioned and auditable rules for fixed services and user-controlled
content; and
\item safe model and policy updates with rollback support.
\end{itemize}

Before browser-store publication, the extension would require a dedicated
security review, permission review, privacy documentation, dependency audit,
and testing on a wider range of Chromium versions and devices. Public release
should also provide a clear statement that the extension is an additional
warning tool and cannot guarantee that a page is safe.

\section{Larger User Evaluation}

A later user study should recruit more participants and use a written test
procedure. Participants could complete a small set of labelled tasks followed
by a period of ordinary browsing. This would permit separate measurement of
task accuracy, warning comprehension, false-warning burden, time to decision,
and user satisfaction.

Consent should remain explicit, and raw browsing data should remain local by
default. If labelled feedback is submitted, the user should see exactly what
will be shared. Repeated paths from the same website should be reported
separately from unique hostnames so that one service does not dominate the
study counts.

\section{Long-Term Maintenance}

Phishing detection is not a one-time training task. A future operational
version would need periodic data-quality checks, source-drift monitoring,
threshold review, generator evaluation, and model rollback. Metrics should be
reported by source and time period as well as in aggregate. Sudden changes in
FPR, FNR, or the score distribution should trigger investigation before a new
model is released.

\section{Chapter Summary}

The most useful next work is to evaluate later labelled URL snapshots, improve
coverage of difficult legitimate and phishing cases, and conduct a larger
supervised browser study. Adding DNS, certificate, reputation, redirect, or
page evidence may further reduce the limitations of URL-only detection.
However, every change should preserve the controlled comparison used in this
project: fixed training conditions, validation-based decisions, untouched test
sets, and transparent reporting of both improvements and failures.
